\section{Exact model potentials from an ample central line}

Retain an original test configuration over an infinite field, with integral
total scheme and dominant parameter projection, an invertible total line
$M$, its original action linearization and its fibre-one identification
with an invertible line $L$. The central fibre may be reducible or nonreduced.
All analytic points and section norms below belong to these original data.

\begin{lemma}[A generated power has zero normalized base-ideal logarithm]
\label{lem:central-ample-generated-power-zero-logarithm}
\uses{prop:generated-line-whole-base-ideal-unit}
\lean{EnvelopeProperty.ProjectiveCentralFiberEmbeddingNeighborhoodClosedWholeGeneratorExactFS.powerBaseIdeal_eq_top_of_generatingSections,EnvelopeProperty.ProjectiveCentralFiberEmbeddingNeighborhoodClosedWholeGeneratorExactFS.gaussPowerBaseIdealLogNormalized_eq_zero_of_generatingSections}
\leanok
If the whole original $M^{\otimes m}$ is generated by global sections,
its all-section base ideal is the unit ideal. Consequently its model norm
is one at every original Gauss point, and its normalized logarithm is zero.
\end{lemma}
\begin{proof}
\leanok
The global generators give an epimorphic evaluation map. Tensor it with
the inverse line and compose with the invertible-line pairing. Its image
is the whole structure sheaf and is contained in the all-section base
ideal. This proves that this ideal is the unit ideal. The valuation of
the unit ideal is one. Taking its logarithm gives zero, and multiplying
by the reciprocal of the degree preserves zero, including degree zero.
\end{proof}

\begin{theorem}[Exact finite rational FS realization of the original potential]
\label{thm:central-ample-original-exact-finite-fs}
\uses{thm:central-embedding-original-centralample-positivepower-whole-generators,
thm:native-tc-original-potential-weighted-fs,
lem:central-ample-generated-power-zero-logarithm}
\lean{EnvelopeProperty.ProjectiveCentralFiberEmbeddingNeighborhoodClosedWholeGeneratorExactFS.finiteGeneratorSections,EnvelopeProperty.ProjectiveCentralFiberEmbeddingNeighborhoodClosedWholeGeneratorExactFS.finiteGeneratorSections_norm,EnvelopeProperty.ProjectiveCentralFiberEmbeddingNeighborhoodClosedWholeGeneratorExactFS.weightedFSData_eval_eq_modelPotential_of_generatingSections,EnvelopeProperty.ProjectiveCentralFiberEmbeddingNeighborhoodClosedWholeGeneratorExactFS.isFinFS_modelPotential_of_central_ample}
\leanok
If the restriction of $M$ to the original central fibre is ample, its
actual native model potential at radius $e^{-1}$ is a finite rational
FS function of the original fibre-one line $L$ and class one.
The same finite family of total-space generators is chosen before all
original analytic points, and its actual Laurent coefficient datum has
exactly that potential as its value at every such point.
\end{theorem}
\begin{proof}
\leanok
Central ampleness supplies a positive integer $m$ and a finite family
of global sections generating $M^{\otimes m}$ on the entire original
total scheme. At an original Gauss point, every generator has model norm
at most one. Pull the generating family back to the valuation centre
and trivialize the line over its local ring. These original coefficients
generate the unit ideal, so at least one coefficient is a unit. Its
valuation is one. Hence the maximum of the same finite generator norms
is exactly one, simultaneously at every original point.

Extract the actual finite Laurent coefficient supports of these sections
using the original action and fibre-one identifications. The existing
all-section model-norm identity says that their normalized weighted
logarithmic maximum is the original model potential plus $m^{-1}$ times
the logarithm of the preceding generator-norm maximum. That maximum is
one, so this correction is zero. Thus the extracted datum evaluates to
the original potential everywhere. This potential is a finite real
number, proving that the extracted original sections have no common
analytic zero. At radius $e^{-1}$ all weights are integer multiples of
$-1$ and are rational. The same datum therefore gives the asserted
finite rational FS function.
\end{proof}

\begin{theorem}[The all-degree asymptotic correction is exactly zero]
\label{thm:central-ample-original-all-degree-zero-correction}
\uses{lem:central-ample-generated-power-zero-logarithm,
thm:central-embedding-original-centralample-positivepower-whole-generators,
thm:native-tc-original-log-limit-and-supremum}
\lean{EnvelopeProperty.ProjectiveCentralFiberEmbeddingNeighborhoodClosedWholeGeneratorExactFS.gaussPowerBaseIdealLogNormalized_tendsto_zero_of_central_ample}
\leanok
Assume also that the original fibre-one line $L$ is ample. At every
original point and every radius $0<r\leq1$, the normalized logarithms
of the original all-section power-base ideals converge to zero, and
zero is their least upper bound over all positive degrees.
\end{theorem}
\begin{proof}
\leanok
The existing all-degree theorem gives a finite nonpositive limit, equal
to the least upper bound over positive degrees. The positive generated
power supplied by central ampleness has normalized logarithm zero by
the first lemma. The least-upper-bound property therefore makes the
limit at least zero. Since it is also nonpositive, it is zero. Substituting
this equality in both the convergence and least-upper-bound statements
gives the result for the unchanged original all-degree sequence.
\end{proof}
